\section{相关工作}

\heading{自动训练并行化}
已有多个系统被提出，用于自动设计深度学习模型训练的并行化策略~\cite{alpa-osdi22,nnscaler-osdi24,unity-osdi22}。这些系统会在 DP、PP 和 TP 维度上进行全尺度模型规划。尽管它们使用穷举搜索算法生成近似最优的训练配置，规划过程往往需要大量时间，通常数分钟才能完成~\cite{nnscaler-osdi24}。LMM 训练需要动态、自适应的流水线调度，因此此类方法并不实用。

\heading{多模态模型训练系统}
研究者已开发许多系统优化~\cite{distmm-nsdi24,optimus-atc25,disttrain-arxiv24,spindle-asplos25,orch-mllm-arxiv25,cornstarch-arxiv25,diffusionpipe-mlsys24}，以应对多模态模型训练中的架构异构性和训练数据异构性。例如，DistMM~\cite{distmm-nsdi24} 处理传统、基于 CLIP 的模型中的模型与数据异构性；在这类模型中，所有编码器并行运行。然而，该设计无法泛化到 LMM，因为 LMM 的编码器、骨干和解码器之间存在顺序依赖。类似地，Spindle~\cite{spindle-asplos25} 和 Optimus~\cite{optimus-atc25} 都以预先定义训练任务的多模态 LLM 为目标，忽略了 LMM 训练流水线的固有动态性。
具体而言，Spindle 针对静态多任务设置而设计，所有任务都必须在训练开始前预先定义。这与现代 LMM 所需的动态、灵活输入处理能力形成对比，而后者正是 \sys{} 关注的重点。

\heading{流水线并行调度}
除 Megatron-LM 的 1F1B 调度外，研究者还提出了多种替代流水线调度。\newline
这些方法包括 Chimera、Hanayo、zero-bubble pipeline 和 GraphPipe~\cite{chimera-sc21,hanayo-sc23,zero-bubble-iclr24,graph-pipe-asplos25}。Chimera~\cite{chimera-sc21} 引入双向流水线以减少流水线气泡；zero-bubble pipeline~\cite{zero-bubble-iclr24} 则将反向计算解耦为输入梯度和权重梯度两个阶段，从而进一步减少气泡。然而，这些方法假定流水线阶段延迟固定，因此仍会受到动态不均衡问题影响。GraphPipe~\cite{graph-pipe-asplos25} 面向包含异构模块的模型，将计算组织为有向无环图（DAG）而非顺序阶段，从而支持更灵活的调度。\sys{} 的流水线调度搜索器可以扩展，以纳入 GraphPipe 一类自定义调度。

\heading{可变序列长度流水线}
若干系统处理单模态 LLM 训练中的文本序列长度变化~\cite{wlb-llm-osdi25,dynapipe-eurosys24,flexpipe-atc25,byte-transformer-ipdps23,sppo-arxiv25}。例如，DynaPipe~\cite{dynapipe-eurosys24} 使用动态规划优化多任务训练的微批次构造。WLB-LLM~\cite{wlb-llm-osdi25} 提出细粒度的逐文档分片策略，以均衡上下文并行组中的工作负载。FlexPipe~\cite{flexpipe-atc25} 通过在线弹性机制动态调整流水线规模，将处理短序列的 GPU 工作进程重新分配去协助处理较长序列。尽管这些方法对文本模型有效，它们无法处理多模态数据。相比之下，我们的方法在单条流水线内支持多种模态。
